\documentclass[10pt,a4paper]{amsart}
\usepackage[margin=19mm]{geometry}
\usepackage{amsmath,amssymb,mathtools}
\usepackage[scr=rsfs,cal=euler]{mathalfa}
\usepackage{xfrac}
\usepackage[unicode,hidelinks,bookmarks=false]{hyperref}
\newcommand{\ie}{i.e.\ }
\newcommand{\cf}{cf.\ }
\newcommand{\resp}{respectively\ }
\newcommand{\Subsection}{\subsection}
\providecommand{\nobreakdash}{\nobreak\hbox{-}}
\setlength{\emergencystretch}{2em}
\multlinegap0pt
\usepackage{tensor}
\usepackage{amssymb}
\newtheorem{thm}{Theorem}
\newtheorem{lem}{Lemma}
\newtheorem{prop}{Proposition}
\newcommand{\A}{{A}}
\newcommand{\CF}{{\mathcal F}}
\newcommand{\Ch}{\mathsf{Ch}}
\newcommand{\Hom}{\operatorname{Hom}}
\newcommand{\Loc}{\mathsf{Loc}}
\newcommand{\R}{{\mathbb{R}}}
\newcommand{\RH}{\operatorname{RH}}
\newcommand{\id}{\operatorname{Id}}
\title{Higher Riemann--Hilbert Correspondence and Descent for Regular Foliations}
\author{Qingyun Zeng}
\date{Source: arXiv:2503.08457v2 (CC BY 4.0)}
\begin{document}
\maketitle

\noindent\emph{Source and licence.} Higher Riemann--Hilbert Correspondence and Descent for Regular Foliations. Version arXiv:2503.08457v2 (CC BY 4.0), distributed by arXiv under the Creative Commons Attribution 4.0 International licence, http://creativecommons.org/licenses/by/4.0/, as recorded in the arXiv licence metadata for this exact version and shown on the abstract page https://arxiv.org/abs/2503.08457v2. Full licence notice: \texttt{licenses/arxiv-2503.08457-CC-BY-4.0.txt}.\par\smallskip

\noindent\textit{Editorial scope.} Counted unit: the article's thm thm:local-rh (source0.tex lines 2412--2434). The article's complete original proof of it is delivered with it (source0.tex lines 2436--2521, one \texttt{proof} environment of 86 lines). No mathematics is written, repaired or re-derived: the statement and the proof are the article's own lines, in the article's own notation, and no retained line is substituted or re-flowed.  Retained uncounted as the source-local context the proof needs: lines 1116--1128; lines 2024--2038; lines 2113--2141; lines 2190--2199; lines 2253--2272; lines 2360--2369.  Nothing was omitted from any retained span, so the package carries no omission record.  No retained line contains a citation, so the wrapper carries no bibliography and no bibliography entry.  No retained line contains a figure environment, an image-inclusion command or drawing code, so no image is delivered and the package declares no asset dependency.  Editorial formatting only, in addition: an AMS \texttt{amsart} preamble that restates the article's own declarations and macros used by the retained lines, namely the environments \texttt{thm}, \texttt{lem}, \texttt{prop} on separate counters, together with the standard packages those lines need.  The retained environments print in this document as Theorem 1, Lemma 1, Proposition 1 in order of appearance, whereas the article prints the same items with the numbering of its own full document.  The complete native source of the article as distributed in the arXiv e-print for this exact version is bundled unchanged as \texttt{sources/174579/source0.tex}.
\begin{thm}[Twisted leafwise de Rham theorem]\label{thm:twisted-dr}
Parallel transport inside each simplex to the first vertex defines a
natural quasi-isomorphism
\[
 I_L:
 \bigl(\Omega^\bullet_{\CF}(M;L),d_\nabla\bigr)
 \longrightarrow
 \bigl(C^\bullet_{\mathrm{sh}}(\CF;L),\delta_\nabla\bigr).
\]
The same conclusion holds for a globally bounded finite-rank graded
flat coefficient bundle with a parallel internal differential, after
totalization.
\end{thm}

\begin{lem}[Leading term of $\RH_1$]\label{lem:rh-leading}
Filter a source morphism by leafwise form degree and a target morphism
by simplicial arity.  Then $\RH_1$ preserves these decreasing
filtrations.  After taking cohomology with respect to the internal
differentials, put
\[
 H_i=H^\bullet(E_i,\mathbb E_i^0),
 \qquad
 L=\underline{\operatorname{Hom}}(H_1,H_2).
\]
The map on the resulting $E_1$-complexes is leafwise integration with
coefficients in $L$ and parallel transport for its induced flat
connection, after the two-ended target complex is identified as in
Lemma~\ref{lem:target-e1}.
\end{lem}

\begin{lem}[The two-ended target $E_1$-complex]
\label{lem:target-e1}
Let $(E_i,\mathbb E_i)$ be regular cohesive modules, let
\[
 H_i=H^\bullet(E_i,\mathbb E_i^0),
 \qquad
 L=\underline{\Hom}(H_1,H_2),
\]
and give $L$ the flat Hom connection.  After taking internal
cohomology, the arity-$p$ associated-graded target term consists of
plot-smooth maps
\[
 a_p(\sigma):
 H_1|_{\sigma(v_p)}
 \longrightarrow
 H_2|_{\sigma(v_0)}.
\]
It is naturally isomorphic, as a sheaf complex, to the one-sided
flat-coefficient cochain complex
$C^\bullet_{\mathrm{sh}}(\CF;L)$.  On a local raw representative the
isomorphism is
\begin{equation}\label{eq:two-ended-xi}
 \Xi_p(a)(\sigma)
 =
 a_p(\sigma)\circ
 \tau^1_{\sigma(v_0)\to\sigma(v_p)}
 \in L_{\sigma(v_0)}.
\end{equation}
\end{lem}

\begin{prop}[Quasi-full faithfulness]\label{prop:qff}
The restriction
\[
 \RH:\mathcal P_{\A}^{\mathrm{reg}}
 \longrightarrow
 \Loc_{\mathrm{shHom},\Ch_{\R}}^{\mathrm{sm,perf,reg}}
 \bigl(\Pi_{\infty}(\CF)\bigr)
\]
is $A_{\infty}$ quasi-fully faithful.
\end{prop}

\begin{lem}[Strictly unital prism transformation]
\label{lem:prism-transformation}
Let $V$ be a strictly unital normalized Block--Smith
infinity-local system on $Y_\bullet$, with structure maps $F_m$.
Define
\begin{equation}\label{eq:unital-prism-components}
 q_0(x)=F_1(h_0x),
 \qquad
 q_n(\sigma)=
 \sum_{j=0}^n(-1)^jF_{n+1}(h_j\sigma)
 \quad(n\geq1).
\end{equation}
Then $q=\{q_n\}$ is a strictly unital normalized closed degree-zero
morphism
\[
 q:g^*V\longrightarrow f^*V.
\]
If $\mathcal H$ is plot-smooth and $V$ is smooth, then $q$ is a raw
plot-smooth morphism.
\end{lem}

\begin{lem}[Raw Whitehead criterion]\label{lem:raw-whitehead}
Let $q:V\to W$ be a closed degree-zero morphism in
$\Loc_{\mathrm{raw}}$.  Suppose that its arity-zero component
\[
 q_0:(V,d_V)\longrightarrow(W,d_W)
\]
has a global smooth chain-homotopy inverse, with global smooth chain
homotopies.  Then $q$ is invertible in
$H^0(\Loc_{\mathrm{raw}})$.
\end{lem}

\begin{thm}[Local higher Riemann--Hilbert correspondence]
\label{thm:local-rh}
Let $B=P\times T$ be a foliated box in which $P$ and $T$ are balls,
$P$ is star-shaped about $p_0$, and the plaques are
$P\times\{t\}$.  Then the raw and fixed-object sheafified-Hom
integration functors
\[
 \RH_{\mathrm{raw}}:
 \mathcal P_{\Omega^\bullet_{\CF}(B)}^{\mathrm{reg}}
 \longrightarrow
 \Loc_{\mathrm{raw},\Ch_{\R}}^{\mathrm{sm,perf,reg}}
 \bigl(\Pi_\infty(\CF|_B)\bigr)
\]
and
\[
 \RH_{\mathrm{shHom}}:
 \mathcal P_{\Omega^\bullet_{\CF}(B)}^{\mathrm{reg}}
 \longrightarrow
 \Loc_{\mathrm{shHom},\Ch_{\R}}^{\mathrm{sm,perf,reg}}
 \bigl(\Pi_\infty(\CF|_B)\bigr)
\]
are $A_\infty$ quasi-equivalences.
\end{thm}

\begin{proof}
Let $r:B\to T$ be projection and $i(t)=(p_0,t)$.  The contraction
\[
 H_s(p,t)=((1-s)p+sp_0,t)
\]
induces a simplicial homotopy from $\id$ to $ir$: for a monotone map
$\alpha:[n]\to[1]$, with affine realization
$|\alpha|:\Delta^n\to I$, put
\[
 \mathcal H_n(\sigma,\alpha)(u)
 =
 H_{|\alpha|(u)}(\sigma(u)).
\]
Every foliated simplex in $B$ has constant transverse coordinate.
It follows that the displayed simplex is again foliated.  Faces and
degeneracies commute with the construction because affine realization
is functorial, and the adjoint formula proves plot-smoothness.

Let $V$ be a target object.  Lemma~\ref{lem:prism-transformation}
gives a closed raw morphism
\begin{equation}\label{eq:local-prism-comparison}
 q_V:r^*i^*V\longrightarrow V,
 \qquad
 (q_V)_0(x)=F_1(h_0x),
 \qquad
 (q_V)_n(\sigma)=
 \sum_{j=0}^n(-1)^jF_{n+1}(h_j\sigma)
 \quad(n\geq1).
\end{equation}
For $x=(p,t)$ its arity-zero component is
\[
 (q_V)_0(x)=F_1(\gamma_x):
 V_{(p_0,t)}\longrightarrow V_x,
 \qquad
 \gamma_x(s)=H_s(x).
\]
Let $\bar\gamma_x(s)=\gamma_x(1-s)$.  The two foliated
$2$-simplices
\begin{align*}
 \Sigma_x(u_0,u_1,u_2)
 &=
 \bigl((u_0+u_2)p+u_1p_0,t\bigr),\\
 \bar\Sigma_x(u_0,u_1,u_2)
 &=
 \bigl((u_0+u_2)p_0+u_1p,t\bigr)
\end{align*}
have boundary paths
\[
 (\gamma_x,\bar\gamma_x,\operatorname{const}_x)
 \quad\text{and}\quad
 (\bar\gamma_x,\gamma_x,\operatorname{const}_{(p_0,t)}),
\]
respectively.  Their Maurer--Cartan equations and strict unitality
show that $F_1(\bar\gamma_x)$ is a two-sided chain-homotopy inverse to
$F_1(\gamma_x)$.  These maps and homotopies vary smoothly with $x$.

Put $K=i^*V$ and let
\[
 E_V=(r^*K,d_K,\nabla^{\mathrm{triv}})
\]
be the pullback complex with its trivial flat leafwise connection and
no higher connection components.  Since $T$ is a ball, the bounded
finite-rank graded bundle $K$ is trivializable, so $E_V$ is cohesive;
regularity is inherited from $V$.  Strict unitality and normalization
give the literal equality
\[
 \RH(E_V)=r^*i^*V.
\]
The unary transport is the identity and all normalized higher
transports vanish.  The construction above gives a global smooth
chain-homotopy inverse to $(q_V)_0$.  Therefore
Lemma~\ref{lem:raw-whitehead} makes
\eqref{eq:local-prism-comparison} an equivalence in the raw target.
Applying the strict dg germ functor $\mathbf{Sh}$ gives an equivalence
in the fixed-object sheafified-Hom target as well.

It remains only to note quasi-full faithfulness of the raw local
functor.  The coefficient contraction in the proof of
Theorem~\ref{thm:twisted-dr} is already a deformation retraction of
the raw coefficient complexes on $B$.  Applying the finite
form-degree/arity filtration and
Lemma~\ref{lem:rh-leading} before sheafification therefore proves that
the raw $\RH_1$ is a quasi-isomorphism.  The sheafified-Hom case is
Proposition~\ref{prop:qff}.  Hence both functors are quasi-fully
faithful and essentially surjective.
\end{proof}

\end{document}
